\setlength{\tabcolsep}{0pt}
\begin{longtable}{|C{12.5mm}|C{12.5mm}|C{12.5mm}|C{12.5mm}|C{12.5mm}|C{12.5mm}|C{12.5mm}|C{12.5mm}|C{12.5mm}|C{12.5mm}|}
\hline
{\bf{HEX}} & {\bf{DEC}} & {\bf{DB7}} & {\bf{DB6}} & {\bf{DB5}} & {\bf{DB4}} & {\bf{DB3}} & {\bf{DB2}} & {\bf{DB1}} & {\bf{DB0}} \\
\hline
\endfirsthead
\multicolumn{3}{l@{}}{\ldots vervolgd}\\
\hline
{\bf{HEX}} & {\bf{DEC}} & {\bf{DB7}} & {\bf{DB6}} & {\bf{DB5}} & {\bf{DB4}} & {\bf{DB3}} & {\bf{DB2}} & {\bf{DB1}} & {\bf{DB0}} \\
\endhead
\multicolumn{3}{l@{}}{vervolgd \ldots}\\
\endfoot
\hline
\endlastfoot
\footnotesize  D080\index{Registers!\$D080} \index{Registers!53376} & \footnotesize 53376 & \footnotesize IRQ & \footnotesize LED & \footnotesize MOTOR & \footnotesize SWAP & \footnotesize SIDE & \multicolumn{3}{c|}{\footnotesize DS\index{Registers!DS}}\\
\hline
\footnotesize  D081\index{Registers!\$D081} \index{Registers!53377} & \footnotesize 53377 & \footnotesize WRCMD & \footnotesize RDCMD & \footnotesize FREE & \footnotesize STEP & \footnotesize DIR & \footnotesize ALGO & \footnotesize ALT & \footnotesize NOBUF \\
\hline
\footnotesize  D082\index{Registers!\$D082} \index{Registers!53378} & \footnotesize 53378 & \footnotesize BUSY & \footnotesize DRQ & \footnotesize EQ & \footnotesize RNF & \footnotesize CRC & \footnotesize LOST & \footnotesize PROT & \footnotesize TK0 \\
\hline
\footnotesize  D083\index{Registers!\$D083} \index{Registers!53379} & \footnotesize 53379 & \footnotesize RDREQ & \footnotesize WTREQ & \footnotesize RUN & \footnotesize WGATE & \footnotesize DISKIN & \footnotesize INDEX & \footnotesize IRQ & \footnotesize DSKCHG \\
\hline
\footnotesize  D084\index{Registers!\$D084} \index{Registers!53380} & \footnotesize 53380 & \multicolumn{8}{c|}{\footnotesize TRACK\index{Registers!TRACK}}\\
\hline
\footnotesize  D085\index{Registers!\$D085} \index{Registers!53381} & \footnotesize 53381 & \multicolumn{8}{c|}{\footnotesize SECTOR\index{Registers!SECTOR}}\\
\hline
\footnotesize  D086\index{Registers!\$D086} \index{Registers!53382} & \footnotesize 53382 & \multicolumn{8}{c|}{\footnotesize SIDE\index{Registers!SIDE}}\\
\hline
\footnotesize  D087\index{Registers!\$D087} \index{Registers!53383} & \footnotesize 53383 & \multicolumn{8}{c|}{\footnotesize DATA\index{Registers!DATA}}\\
\hline
\footnotesize  D088\index{Registers!\$D088} \index{Registers!53384} & \footnotesize 53384 & \multicolumn{8}{c|}{\footnotesize CLOCK\index{Registers!CLOCK}}\\
\hline
\footnotesize  D089\index{Registers!\$D089} \index{Registers!53385} & \footnotesize 53385 & \multicolumn{8}{c|}{\footnotesize STEP\index{Registers!STEP}}\\
\hline
\footnotesize  D08A\index{Registers!\$D08A} \index{Registers!53386} & \footnotesize 53386 & \multicolumn{8}{c|}{\footnotesize PCODE\index{Registers!PCODE}}\\
\hline
\end{longtable}
\begin{itemize}
\item{\bf{ALGO}\index{Registers!ALGO}} Selecteert het lees- en schrijfalgoritme (wordt momenteel genegeerd).
\item{\bf{ALT}\index{Registers!ALT}} Selecteert de alternatieve DPLL-leesherstelmethode (niet geïmplementeerd)
\item{\bf{BUSY}\index{Registers!BUSY}} F011 FDC-busyvlag (commando wordt uitgevoerd) (alleen lezen)
\item{\bf{CLOCK}\index{Registers!CLOCK}} Stelt het klokpatroon in dat wordt gebruikt bij het schrijven van adres- en datamerken. Normaal \$FF laten staan
\item{\bf{COMMAND}\index{Registers!COMMAND}} F011 FDC-commandoregister
\item{\bf{CRC}\index{Registers!CRC}} F011 FDC CRC-controlefoutvlag (alleen lezen)
\item{\bf{DATA}\index{Registers!DATA}} F011 FDC-gegevensregister (lezen/schrijven) voor toegang tot de 512 byte sectorbuffer van de schijfstationcontroller
\item{\bf{DIR}\index{Registers!DIR}} Stelt de staprichting in (naar binnen ten opzichte van
\item{\bf{DISKIN}\index{Registers!DISKIN}} F011 Schijfgevoeligheid (alleen lezen)
\item{\bf{DRQ}\index{Registers!DRQ}} F011 FDC DRQ-vlag (een of meer bytes gegevens zijn gereed) (alleen lezen)
\item{\bf{DS}\index{Registers!DS}} Stationselectie (0 tot 7). Intern station is 0. Tweede schijfstation op de interne kabel is 1. Andere waarden zijn gereserveerd voor de C1565 externe stationinterface.
\item{\bf{DSKCHG}\index{Registers!DSKCHG}} F011 Schijfwisselgevoeligheid (alleen lezen)
\item{\bf{EQ}\index{Registers!EQ}} F011 FDC CPU- en diskschijfwaarden naar de sectorbuffer zijn gelijk, wat aangeeft dat de sectorbuffer vol of leeg is. (alleen lezen)
\item{\bf{FREE}\index{Registers!FREE}} Commando is een vrij formaat (laag niveau) bewerking
\item{\bf{INDEX}\index{Registers!INDEX}} F011 Indexgatgevoeligheid (alleen lezen)
\item{\bf{IRQ}\index{Registers!IRQ}} De schijfstationcontroller heeft een interrupt gegenereerd (alleen lezen). Opmerking: interrupts worden momenteel niet geïmplementeerd op de 45IO27.
\item{\bf{LED}\index{Registers!LED}} Station-LED knippert wanneer ingesteld
\item{\bf{LOST}\index{Registers!LOST}} F011 LOST-vlag (gegevens zijn verloren gegaan tijdens de overdracht, d.w.z. de CPU las de gegevens niet snel genoeg) (alleen lezen)
\item{\bf{MOTOR}\index{Registers!MOTOR}} Schakelt de stationmotor en LED in (tenzij het LED-signaal ook is ingesteld, waardoor de station-LED knippert)
\item{\bf{NOBUF}\index{Registers!NOBUF}} Herstel de lees-/schrijfpunten van de sectorbuffer
\item{\bf{PCODE}\index{Registers!PCODE}} (Alleen lezen) geeft de beschermingscode van het meest recent gelezen sector terug. Was bedoeld voor elementaire kopieerbescherming. Niet geïmplementeerd.
\item{\bf{PROT}\index{Registers!PROT}} F011 Schrijfbeveiligingsvlag van de schijf (alleen lezen)
\item{\bf{RDCMD}\index{Registers!RDCMD}} Commando is een leesbewerking wanneer ingesteld
\item{\bf{RDREQ}\index{Registers!RDREQ}} F011 Leesverzoekvlag, d.w.z. het gevraagde sector werd gevonden tijdens een leesbewerking (alleen lezen)
\item{\bf{RNF}\index{Registers!RNF}} F011 FDC-verzoek niet gevonden (RNF), d.w.z. een lees- of schrijfbewerking heeft het gevraagde sector niet gevonden (alleen lezen)
\item{\bf{RUN}\index{Registers!RUN}} F011 Opeenvolgende overeenkomst. Een synoniem van RDREQ op de 45IO27 (alleen lezen)
\item{\bf{SECTOR}\index{Registers!SECTOR}} F011 FDC-sectorselectieregister
\item{\bf{SIDE}\index{Registers!SIDE}} Stuurt het SIDE-signaal naar de schijfstationcontroller direct, d.w.z. selecteert welke zijde van het medium actief is.
\item{\bf{STEP}\index{Registers!STEP}} Het schrijven van 1 laat de kop in de aangegeven richting stappen
\item{\bf{SWAP}\index{Registers!SWAP}} Wissel de boven- en onderste helft van de gegevensbuffer (d.w.z. inverteer bit 8 van de sectorbuffer)
\item{\bf{TK0}\index{Registers!TK0}} F011 Kop bevindt zich boven spoor 0-vlag (alleen lezen)
\item{\bf{TRACK}\index{Registers!TRACK}} F011 FDC-spoorselectieregister
\item{\bf{WGATE}\index{Registers!WGATE}} F011 Schrijfpoortvlag. Geeft aan dat het station momenteel naar het medium schrijft. Er kunnen problemen optreden als een schrijfbewerking wordt afgebroken (alleen lezen)
\item{\bf{WRCMD}\index{Registers!WRCMD}} Commando is een schrijfbewerking wanneer ingesteld
\item{\bf{WTREQ}\index{Registers!WTREQ}} F011 Schrijfverzoekvlag, d.w.z. het gevraagde sector werd gevonden tijdens een schrijfbewerking (alleen lezen)
\end{itemize}
